\section{An ultrafilter rank and simultaneous reservations}
\label{sec:filters}

We construct two tools on a countable set. The first is an ordinal-valued
function on an ultrafilter whose value can be increased while preserving
a prescribed lower bound on the set. The second supplies two disjoint
sets that meet each of fewer than continuum many prescribed sets in
positive upper block density. These two tools will allow us to choose bases
with an ordinal inequality while also separating them from all earlier
choices.

\subsection{Blocks, columns, and two ideals}

Let $\mathfrak c=2^{\aleph_0}$, also regarded as its initial ordinal.
For each integer $m\ge0$, let
\[
 D_m=\{m\}\times\{0,1\}^{\{0,1\}^m},\qquad
 D=\coprod_{m\ge0}D_m,\qquad w_m=|D_m|=2^{2^m}.
\]
Thus $D$ is countably infinite. The tags $m$ make its finite blocks
disjoint. Their sizes satisfy
\begin{equation}\label{eq:block-growth}
 w_m\ge2\sum_{j<m}w_j\qquad(m\ge1).
\end{equation}
Indeed, equality holds for $m=1$. If the inequality holds for $m$, then
$2\sum_{j<m+1}w_j\le3w_m\le w_m^2=w_{m+1}$, since $w_m\ge4$.

The density-zero ideal on $D$ is
\[
 \mathcal J=\left\{X\subseteq D:
       \frac{|X\cap D_m|}{w_m}\longrightarrow0\right\}.
\]
An ideal is a family closed under taking subsets and finite unions;
we call a set \emph{positive} for the ideal if it does not belong to it.
The displayed family is a proper ideal containing all finite sets.

Let $L$ and $R$ be labelled copies of $D$, and write
$E_0=L\mathbin{\dot\cup}R$. For $X\subseteq E_0$, its two slices,
viewed as subsets of $D$, are denoted by $X_L$ and $X_R$.
When writing lifts, we identify $E_0$ with $D\times\{L,R\}$.
Write $W_m$ for the union of the two copies of $D_m$.
The ideal $\mathcal J_0$ on $E_0$ consists of the sets whose two slices
belong to $\mathcal J$. Equivalently,
\[
 X\in\mathcal J_0
 \quad\Longleftrightarrow\quad
 \frac{|X\cap W_m|}{|W_m|}\longrightarrow0.
\]
This equivalence follows because the doubled-block density is the
average of the two nonnegative slice densities. Complements are always
taken in the relevant ground set, $D$ or $E_0$.

\begin{lemma}[Independent columns]\label{lem:columns}
There are subsets $A_t,B_t,G_t$ of $D$, indexed by $t<\mathfrak c$,
with the following property. For any $r$ distinct columns from these
three families and any prescribed membership pattern on them, the
proportion of that pattern in $D_m$ is exactly $2^{-r}$ for all
sufficiently large $m$.
\end{lemma}

\begin{proof}
Assign distinct infinite binary strings to all the columns; this is
possible because $3\mathfrak c=\mathfrak c$. For a column with string
$h$, put $(m,f)$ in that column precisely when
$f(h\upharpoonright m)=1$.
The prefixes of any fixed finite list of distinct strings are distinct
for all sufficiently large $m$. Prescribing the values of $f$ at these
$r$ arguments leaves $2^{2^m-r}$ functions among the $2^{2^m}$ choices
in $D_m$, proving the assertion.
\end{proof}

We need a larger ideal in addition to $\mathcal J$. Define
\begin{equation}\label{eq:larger-ideal}
 \mathcal K=\left\{X\subseteq D:
 X\cap\bigcap_{t\in F}A_t\in\mathcal J
 \text{ for some finite }F\subseteq\mathfrak c\right\},
\end{equation}
where an empty intersection is $D$. This is the ideal generated by
$\mathcal J$ and the complements $A_t^c$. To see the equivalence, split
$X$ into its intersection with the positive $A$-columns in
\eqref{eq:larger-ideal} and its remaining part, which is covered by
their complements. In particular $\mathcal J\subseteq\mathcal K$.
Lemma~\ref{lem:columns} shows that $D\notin\mathcal K$, so this ideal
is proper.

We will repeatedly use the following consequence. A consistent finite
pattern prescribing positive membership in $A$-columns and arbitrary
signs in $G$-columns is $\mathcal K$-positive. Intersecting such a pattern
with any further finitely many positive $A$-columns still gives a
pattern with a fixed positive eventual density, so
\eqref{eq:larger-ideal} excludes membership in $\mathcal K$.

Write $X\subseteq_{\mathcal K}Y$ if $X\setminus Y\in\mathcal K$.
On $E_0$, let $\mathcal K_0$ be the ideal of sets whose two slices
belong to $\mathcal K$. Thus $\mathcal J_0\subseteq\mathcal K_0$.
The distinction between these ideals matters: changes in
$\mathcal K_0$ will preserve an ordinal condition, whereas differences
positive for $\mathcal J_0$ will keep the resulting bases incomparable.

\subsection{The ultrafilter and its ordinal rank}

An ultrafilter $\mathcal U$ on $D$ is a proper filter that contains
exactly one of $X$ and $D\setminus X$ for every $X\subseteq D$.

\begin{lemma}\label{lem:ultrafilter}
There is an ultrafilter $\mathcal U$ on $D$ disjoint from $\mathcal K$,
containing every $G_t$, and satisfying
\begin{equation}\label{eq:finitely-many-columns}
 X\in\mathcal U\quad\Longrightarrow\quad
 \{t<\mathfrak c:X\subseteq_{\mathcal K}G_t\}
 \text{ is finite}.
\end{equation}
\end{lemma}

\begin{proof}
Call $Z\subseteq D$ forbidden if it is contained modulo $\mathcal K$
in infinitely many of the $G_t$. Require a filter to contain all
complements of members of $\mathcal K$, all $G_t$, and all complements
of forbidden sets. We verify the finite-intersection property.

Combine the finitely many $\mathcal K$-small exclusions into one set
$N\in\mathcal K$. Let $F$ be the finite set of required positive
$G$-indices, and let $Z_1,\ldots,Z_q$ be the forbidden sets to exclude.
For each $j$, choose an index $t_j$ such that
$Z_j\setminus G_{t_j}\in\mathcal K$, with the $t_j$ distinct and
outside $F$. Each $Z_j$ has infinitely many witnessing indices, so
these choices are possible. The pattern
\[
 P=\bigcap_{t\in F}G_t\cap\bigcap_{j=1}^qG_{t_j}^c
\]
is $\mathcal K$-positive. Moreover
$P\cap Z_j\subseteq Z_j\setminus G_{t_j}\in\mathcal K$.
It follows that
\[
 P\setminus\left(N\cup\bigcup_{j=1}^qZ_j\right)
\]
is still $\mathcal K$-positive: otherwise $P$ would be a finite union
of members of $\mathcal K$. This nonempty set lies in the required
finite intersection. Extend the resulting proper filter to an
ultrafilter by the ultrafilter lemma. It contains no member of
$\mathcal K$ and no forbidden set, as required.
\end{proof}

Fix such an ultrafilter. Call its members \emph{large} and its
nonmembers \emph{small}. Small in this sense does not mean membership
in either ideal. Every large set is $\mathcal J$-positive, and every
complement of a $\mathcal K$-set is large.

List all large subsets of $D$ as $(F_\xi)_{\xi<\mathfrak c}$,
allowing repetitions if needed. For a large set $X$, define
\begin{equation}\label{eq:ordinal-rank}
 \rho(X)=\min\{\xi<\mathfrak c:
                     F_\xi\subseteq_{\mathcal K}X\}.
\end{equation}
The set of indices is nonempty because $X$ occurs in the list.
This is an auxiliary ordinal rank, distinct from matroid rank.

\begin{lemma}[Rank invariance and monotonicity]\label{lem:rank}
If $X\triangle Y\in\mathcal K$, then $X$ is large if and only if $Y$
is large, and in that case $\rho(X)=\rho(Y)$.
If $X,Y$ are large and $X\subseteq_{\mathcal K}Y$, then
$\rho(X)\ge\rho(Y)$.
\end{lemma}

\begin{proof}
The complement of $X\triangle Y$ belongs to $\mathcal U$.
Intersecting it with either large set gives a large subset of the
other, proving invariance of largeness. The two sets have exactly the
same candidate indices in \eqref{eq:ordinal-rank}, proving rank
invariance. Finally, $F_\xi\subseteq_{\mathcal K}X$ and
$X\subseteq_{\mathcal K}Y$ imply
$F_\xi\subseteq_{\mathcal K}Y$ by closure of the ideal under finite
unions. Taking minima proves the last assertion.
\end{proof}

The next property supplies large sets of arbitrarily high rank inside
any interval whose lower endpoint is small and upper endpoint is
large.

\begin{lemma}[Rank-raising thinning]\label{lem:thinning}
Let $s\subseteq X\subseteq D$, with $s$ small and $X$ large, and let
$\gamma<\mathfrak c$. There is a large set $u$ such that
\[
 s\subseteq u\subseteq X,\qquad \rho(u)>\gamma.
\]
\end{lemma}

\begin{proof}
For every $\xi\le\gamma$, the set
$Y_\xi=F_\xi\cap(X\setminus s)$ is large. By
\eqref{eq:finitely-many-columns}, it is contained modulo
$\mathcal K$ in only finitely many $G_t$.
The union of these finite sets of indices has cardinal at most
$\max(\aleph_0,|\gamma+1|)<\mathfrak c$. Choose $t$ outside that union
and put
\[
 u=s\cup(X\cap G_t).
\]
Then $u$ is large and lies between $s$ and $X$.
If $F_\xi\subseteq_{\mathcal K}u$ for some $\xi\le\gamma$, then
\[
 Y_\xi\setminus G_t
 =F_\xi\cap X\cap s^c\cap G_t^c
 \subseteq F_\xi\setminus u\in\mathcal K,
\]
contrary to the choice of $t$. Thus $\rho(u)>\gamma$.
\end{proof}

This counting argument does not require regularity of
$\mathfrak c$. For any fixed cardinal $\lambda<\mathfrak c$, a union
of $\lambda$ finite sets has cardinal at most
$\max(\lambda,\aleph_0)<\mathfrak c$. We will use this same estimate
at each stage of a recursion of length $\mathfrak c$.

\subsection{Two simultaneous reservations}

The ordinal rank is unchanged by $\mathcal K$-small changes. We now
show that such changes can nevertheless meet many prescribed sets
in a way that remains positive for the smaller ideal $\mathcal J$.
For $t<\mathfrak c$, define the following subsets of $E_0$, using the
same labels in both copies:
\[
 H_{t,0}=(A_t^c\cap B_t^c)\times\{L,R\},\qquad
 H_{t,1}=(A_t^c\cap B_t)\times\{L,R\}.
\]
They are disjoint, and both belong to $\mathcal K_0$ because their
slices lie in $A_t^c$.

\begin{lemma}[Simultaneous reservations]\label{lem:reservations}
Let $\mathscr S$ be a family of fewer than $\mathfrak c$ subsets of
$E_0$, each positive for $\mathcal J_0$. There is an index
$t<\mathfrak c$ such that
\[
 S\cap H_{t,0}\notin\mathcal J_0
 \quad\text{and}\quad
 S\cap H_{t,1}\notin\mathcal J_0
 \qquad(S\in\mathscr S).
\]
\end{lemma}

\begin{proof}
Fix first one $\mathcal J_0$-positive set $S$. Choose $\epsilon>0$
such that the proportion of $S$ in $W_m$ is at least $\epsilon$ for
infinitely many $m$. Call an index $t$ bad for $S$ if at least one
of $S\cap H_{t,0}$ and $S\cap H_{t,1}$ belongs to $\mathcal J_0$.
We prove that only finitely many indices are bad.

Choose any finite list of distinct bad indices $t_1,\ldots,t_q$, and
for each $i$ choose $k_i\in\{0,1\}$ with
$S\cap H_{t_i,k_i}\in\mathcal J_0$. Put $H_i=H_{t_i,k_i}$.
These cells use $2q$ distinct $A$- and $B$-columns. Avoiding $H_i$
allows three of the four membership patterns on its two columns.
Lemma~\ref{lem:columns} therefore gives $3^q$ avoiding patterns out of
$4^q$ equally frequent patterns in every sufficiently late block.
Duplicating the cells on the two copies does not change these proportions.
In particular,
\[
 \frac{|W_m\setminus\bigcup_{i=1}^qH_i|}{|W_m|}
 =\left(\frac34\right)^q
\]
for all sufficiently large $m$. Every point of $S\cap W_m$ either
avoids all these cells or lies in one of the failing intersections, so
\[
 \frac{|S\cap W_m|}{|W_m|}
 \le \left(\frac34\right)^q+
       \sum_{i=1}^q\frac{|S\cap H_i\cap W_m|}{|W_m|}.
\]
Each summand tends to zero. Along the infinite subsequence on which
$S$ has density at least $\epsilon$, this gives
$\epsilon\le(3/4)^q$. Since $(3/4)^q\to0$, the size of any finite
list of bad indices is bounded. Hence only finitely many indices are bad.

For the family $\mathscr S$, take the union of these finite sets of
bad indices. Its cardinal is less than $\mathfrak c$ by the estimate
following Lemma~\ref{lem:thinning}. Any index outside that union
works simultaneously for every $S\in\mathscr S$. If the family is
empty, any index works.
\end{proof}
